\section{The Relay and Custody Signing}\label{sec:relay}

Agents that hold no keys submit Soroban calls through a relay
(\file{crates/grant-relay}) that speaks the protocol Nido's browser client
already uses: the client posts a host function and its signed authorization
entries to \code{/relay} and polls for the result, so a deployment changes only
the client's relay URL\claim{relay-frontend}. The relay runs the native build
of the same evaluator. Its grant bounds what the custody key will sign, pay
for and submit; authority over an account is still decided on chain, by the
account's own \code{\_\_check\_auth} and, for agents, its grant signer. The
repository's other relay, \file{infra/relayer}, wraps a third-party
fee-sponsoring plugin with no grant check and is outside this design: it
relays any contract call whose on-chain authorization is valid, which for an
account with a grant signer still means the grant signer's check. Each agent has a chain grant over the contexts of the
call's authorization trees and a signing grant over the custody signature.
Before any key is used, the relay requires, in order: that the API key names
an agent; that no supplied entry uses the transaction source's credentials;
that the chain grant, as the rule-eval rule set its caveats lower to, allows
the host function and every context of every supplied tree, read through
rule-core's Soroban facade as the grant signer reads them; and that the
signing grant allows the custody key, the
network and the \code{transaction} payload kind. These checks need no network,
so a denial leases nothing and never reaches the key
keeper\claim{relay-nested}.

It then leases a channel account and runs two simulations. A recording
simulation of the bare host function must succeed, must not ask for the
channel's own authorization\claim{relay-channel-auth}, and every context it
records must pass the chain grant\claim{relay-recorded}; an enforcing
simulation with the supplied entries must succeed, and the fee, the simulated
minimum plus a margin for rent that grows before the transaction lands, must
stay under a cap\claim{relay-fee-margin}. The supplied trees are what the
user's signatures authorize, and the recorded trees are what the execution asks for, so a call is relayed only when
the grant allows every context in either. That the two together cover every
context the signatures authorize rests on the host's matching of
\code{require\_auth} to signed trees and is not proven\claim{trees-cover}. The
chain grant does not see calls that ask for no authorization, calls a contract
makes inside its own \code{\_\_check\_auth}, or the fee, which the cap bounds.
Soroban RPC's recording simulation records authorization only for the root
invocation, so the relay cannot submit a call whose first request for an
account's authorization is made in a nested call\claim{network-root-recording}.

\paragraph{Relay grants and chain grants.} The relay's grants are a
configuration file with no link to the grant signer, so a deployment must
rewrite it when the admin document issues, revokes or rotates grants on chain;
that a deployment keeps the two in step is unproven\claim{relay-grants-in-step}.
A relay grant wider than the chain's lets more calls reach the network's
simulation and the grant signer, which still refuse them: with relay grants
allowing every call, all 44 attempts the network runs are
refused\claim{network-attack-suite}, each but one by the component that refuses
it in the simulation\claim{network-as-simulation}. A relay grant narrower than
the chain's refuses calls the account would allow.

\paragraph{Channels and the key keeper.} Transactions go out from a pool of
channel accounts, each with master weight 0 and the custody key as its only
signer; the relay checks this for every channel at start. A request holds its
channel until its transaction settles, so concurrent requests use distinct
channels, the relay reads each channel's sequence number fresh, and each
channel carries at most one transaction per ledger\claim{relay-throughput}.
The relay has no budget, so concurrent spends that together overdraw an
ancestor all pass it; the grant signer refuses the excess on chain, and the
refused transactions cost their channels a fee\claim{network-race}. The key
keeper signs the transaction hash with the custody key and the relay verifies
the signature before attaching it. The keeper is an interface with three
implementations: an in-process ed25519 keeper, a client for the Dfns Keys
API's Generate Signature endpoint with Dfns's user-action signing flow, and a
local mock of that endpoint that follows Dfns's OpenAPI document. Dfns's own
policies remain a check behind the grant, but they see only a 32-byte hash,
which their amount and recipient rules cannot evaluate~\cite{dfnsSigningPolicies}.

\paragraph{Channel accounts on testnet.} A channel's source account is not the
custody key's own address, so the arrangement depends on the network accepting
a transaction whose only signature is a signer's that is not the source. On
Stellar testnet, with the custody key held by the Dfns mock, channel-sourced
classic payments, Soroban invocations and fee bumps signed only by the custody
key were accepted, and the same transactions signed by a channel's zeroed
master key or by an unrelated key were refused with
\code{txBAD\_AUTH}\claim{testnet-channel}. Whether real Dfns signs such
transactions is open: its reference does not say whether the
\code{Transaction} signature kind checks the envelope's source account against
the key's own address\claim{dfns-transaction-source}, or whether a
\code{Hash} request on an ed25519 key signs the 32 bytes directly as the
message\claim{dfns-hash}.
